%% exemplo-relatorio.tex
%% Documento mínimo de teste — linha Relatório A4 CRP-SP
%% Exercita: seccionamento, boxes paralelos, tabela com linha
%% destacada, box de destaque, listas e figura com alt.
%% Compilar 2x com: sh lualatex-podman.sh exemplo-relatorio.tex

\DocumentMetadata{
  lang        = pt-BR,
  pdfversion  = 2.0,
  pdfstandard = ua-2,
  tagging     = on,
}

\documentclass{crpsp-relatorio}

% --- Paleta desta publicação (CRP-SP) ---
\crpspCorPrincipal{422C73}   % Her Highness      — 11,4:1 sobre branco
\crpspCorSecundaria{0511F2}  % Midsummer Nights  —  8,9:1
\crpspCorAcessoria{A55DA6}   % Royal Pretender   —  4,45:1, só gráfico
\crpspCorAlarme{EE05F2}      % Piquant Pink      —  3,54:1, só moldura

\hypersetup{
  pdftitle  = {Exemplo de Relatório CRP-SP},
  pdfauthor = {Conselho Regional de Psicologia de São Paulo},
}

\title{Exemplo de Relatório CRP-SP}
\author{}
\date{}

\begin{document}

\section{A decisão}

Parágrafo de abertura, para dar contexto de fluxo ao tagging e
verificar a medida de linha resultante das margens de 25\,mm.
O corpo do texto vai em preto sobre branco, contraste de 19,4:1.

\subsection{Um subtítulo}

Texto sob o subtítulo, conferindo o espaçamento entre níveis.

\section{As duas opções}

\begin{paralelos}
  \boxlado{Opção A}{Manter o modelo atual}{%
    \begin{itemize}
      \item Periodicidade segue abaixo da meta regimental.
      \item A checagem de fatos continua sem responsável.
      \item O risco de erro material permanece.
    \end{itemize}}
  \boxlado[destaque]{Opção B $\cdot$ recomendada}{Nova divisão de trabalho}{%
    \begin{itemize}
      \item Pauta definida pela ComCom com a Comunicação.
      \item Checagem de fatos atribuída à Comunicação.
      \item Validação final pelo Plenário.
    \end{itemize}}
\end{paralelos}

Parágrafo após os boxes, para confirmar que o fluxo normal foi
retomado e que o espaçamento vertical fecha.

\subsection{Divisão de trabalho proposta}

\begin{tabular}{p{0.4\linewidth}p{0.5\linewidth}}
  \linhacab \celcab{Etapa} & \celcab{Responsável na Opção B} \\
  Definição de pauta & ComCom + Comunicação \\
  Apuração e entrevistas & Conselheiras/os \\
  Redação das matérias & Conselheiras/os \\
  Revisão e padronização & Comunicação \\
  \linhadestaque \textbf{Checagem de fatos} & \textbf{Comunicação} \\
  Diagramação e projeto gráfico & Comunicação \\
  Produção gráfica e distribuição & Comunicação \\
  Validação final & Plenário \\
\end{tabular}

A linha destacada acima é a etapa que hoje não tem responsável
designado. Numa tabela de funções, a célula vazia recebe \vago.

\section{Diagnóstico}

\begin{destaque}{Caso documentado $\cdot$ 2026}
  Em 2026 --- ano com a menor equipe nominal da série, 2 pessoas ---
  a edição 209 foi impressa e postada com dois erros materiais no
  expediente.

  \textbf{Nenhuma etapa do processo era capaz de interceptar o erro
  antes da impressão, porque essa etapa não existe no fluxo atual.}
\end{destaque}

Parágrafo final, fechando o documento de teste.

\end{document}
